%! app: TODOapp
%! outcome: TODOoutcome

An element for which $P(x) = F$ is called a {\bf counterexample} of $\forall x P(x)$. When there is a counterexample in the domain of predicate $P$,  $\forall x P(x)$ is false.


An element for which $P(x) = T$ is called a {\bf witness} of $\exists x P(x)$. When there is a witness in the domain of predicate $P$, then $\exists x P(x)$ is true.
